\section{Hardwarearchitektur}
Die Hardwarearchitektur der WordClock ist modular aufgebaut.
Zentrale Komponente ist ein \gls{esp32c6}-Mikrocontroller, der die Ansteuerung der \gls{led}-Stränge, die Verarbeitung der Sensordaten sowie die drahtlose Kommunikation übernimmt.
Die Energieversorgung erfolgt über ein 5-V-Netzteil mit einer maximalen Ausgangsstromstärke von 10 A.
Über diese gemeinsame Versorgung werden sowohl der Mikrocontroller als auch die drei \gls{ws2812b}-\gls{led}-Stränge betrieben.

Die Anzeige besteht aus drei funktional getrennten \gls{led}-Bereichen.
Die Hauptanzeige wird durch eine 16 $\times$ 16-\gls{led}-Matrix mit insgesamt 256 \gls{led}s gebildet.
Sie dient zur Darstellung der Uhrzeit sowie zusätzlicher Wetter- und Temperaturbegriffe.
Ein separater \gls{led}-Strang mit 60 \gls{led}s wird für die Sekundenanzeige verwendet.
Ein weiterer \gls{led}-Strang mit 31 \gls{led}s bildet die Zusatzinformationen: Innenraumtemperatur, Luftfeuchtigkeit, Tageszeit und Mondphase ab.

Als Sensor wird ein \gls{dht22} eingesetzt.
Dieser erfasst die Innenraumtemperatur und die relative Luftfeuchtigkeit.
Die Messwerte werden vom \gls{esp32c6} eingelesen und anschließend auf dem dritten \gls{led}-Strang mit den Zusatzinformationen dargestellt.

Der mechanische Aufbau basiert auf einer Holzgrundplatte, auf der die \gls{led}-Stränge montiert sind.
Die Front besteht aus 16 miteinander verschraubten 3D-gedruckten Modulen.
Jedes Modul enthält geschlossene Lichtkammern für die einzelnen \gls{led}s.
Dadurch wird eine optische Trennung der Buchstaben erreicht.
Zusätzlich wird pro \gls{led}-Kammer ein weißes, 3D-gedrucktes Diffusorplättchen eingesetzt, um die Lichtverteilung zu homogenisieren.